\documentclass[12pt,a4paper]{article}
\usepackage[spanish=nohyphenation]{hyphsubst}
\usepackage[spanish]{babel}
\usepackage{selinput}
\usepackage{tikz}
\usepackage{blindtext}
\usepackage{csquotes}
\usepackage{graphicx}
\usepackage[font=footnotesize,labelfont=bf,justification=centering,textfont=it]{caption}
\usepackage{float}
\usepackage{hyperref}
\usepackage{tabularx}
\usepackage{longtable}
\usepackage{listings}
\usepackage{color}
\usepackage{amsmath}
\usepackage{mathtools}
\usepackage{datetime}
\usepackage{listingsutf8}
\usepackage{mwe}
\usepackage{biblatex}

\addbibresource{teg/references.bib}

\makeatletter

\setlength\bibitemsep{0.5\baselineskip}

\definecolor{lightgray}{rgb}{.9,.9,.9}
\definecolor{darkgray}{rgb}{.4,.4,.4}
\definecolor{purple}{rgb}{0.65, 0.12, 0.82}

\newdateformat{mydate}{\shortmonthname[\THEMONTH], \THEYEAR}

\renewcommand*{\thefootnote}{\arabic{footnote}}

\graphicspath{ {images/} }
\title{
	{Desarrollo de un editor de visualizaciones de propiedades de historiales de wikis}\\
	{\large Universidad Central de Venezuela}\\
}
\author{Adrian J. Mejias. O.}



\addto\captionsspanish{
  \renewcommand{\listfigurename}
    {índice de figuras}
  \renewcommand{\listtablename}
    {índice de tablas}
  \renewcommand\lstlistingname{Código Fuente}
  \renewcommand\lstlistlistingname{índice de código fuente}
}
\renewcommand\spanishtablename{Tabla}
\newcommand{\tabitem}{~~\llap{\textbullet}~~}

\begin{document}

  \begin{titlepage}
	\centering
	{\large República Bolivariana de Venezuela \par}
	{\large Universidad Central de Venezuela\par}
	{\large Facultad de Ciencias\par}
	{\large Escuela de Computación\par}\vspace{2cm}

	\includegraphics[scale=0.60]{ucv_logo.png}\par\vspace{1cm}
	{\scshape\Large\textbf{\@title}\par}
	\vfill

    {\large Tutor Prof. Eugenio Scalise \par}
    {\large Pasante subpagado \@author \par}
	{\large \mydate\today \par}
\end{titlepage}

  \section*{Resumen}

% Tıtulo:
% Desarrollo de un editor de visualizaciones de propiedades de historiales de
% wikis.
% Autor:
% Leonardo Testa.
% Tutor:
% Prof. Eugenio Scalise.
% Un Wiki es un sitio web, generalmente de car ́acter informativo (como lo es
% Wikipedia), que puede ser modificado por m ́ultiples personas. Cada una de
% estas modificaciones son almacenadas, y en conjunto conforman un historial
% de versiones, en donde cada versi ́on representa una modificaci ́on y los efectos
% que caus ́o en el art ́ıculo wiki. Siendo Wikipedia un caso real con bastante

% popularidad, es normal que el historial de versiones de un art ́ıculo sea su-
% ficientemente extenso y complejo, por lo tanto las personas interesadas en

% mantener el art ́ıculo “sano” perder ́an una gran suma de tiempo revisando
% las modificaciones. En este documento, presentaremos la investigaci ́on y la
% realizaci ́on de una herramienta web que facilita la lectura de propiedades del
% historial a aquellas personas interesadas, en donde se optar ́a por visualizaci ́on

% de datos como estrategia, de esta forma, mediante una interfaz capaz de ma-
% nipular gr ́aficas el usuario podr ́a proyectar distintas propiedades y conseguir

% f ́acilmente informaci ́on m ́as completa y concretar patrones.
% Palabras claves:
% Visualizaci ́on de datos, wiki, propiedades de historiales, gr ́aficas, herramienta
% web, editor de visualizaciones, wikipedia.





  \tableofcontents

  \listoffigures

  \section{Introduccion}

\subsection{Objetivo general}

\subsection{Objetivos específicos}

\subsection{Justificación}

\subsection{Distribución del documento}



  \section{Marco teórico}

  \section{Marco Técnologico}

  \section{Marco Aplicativo}

Revision de metodologias agiles vs pesadas, usos comunes y puntos a favor para seleccionar agile

agiles es mejor, comparemos entre ellas
\cite{ComparacionFrameworks}
\subsection{Revisión de metodologías ágiles para web }

    \subsubsection{Lean}
    \subsubsection{XP}
    \subsubsection{Scrum}
    \subsubsection{TDD}

  \section{Conclusiones}

  \printbibliography

\end{document}